\SourceLeaf{451}
\par\Needspace{8\baselineskip}\bigskip
\pdfbookmark[3]{Exposition, 4}{pallas-II-article4}
\begin{center}4.\end{center}
\noindent\hspace*{\parindent}Il nous reste encore de développer la liaison entre nos constantes et les élémens dont les astronomes ont coutume de faire usage. Pour cet effet, rapportons à une surface sphérique d’un rayon arbitraire et dont le centre coïncide avec celui du soleil, les neuf directions suivantes: les trois axes des coordonnées; la droite menée du soleil vers la planète; une droite perpendiculaire à celle-ci dans le plan de l’orbite; la droite perpendiculaire à ce plan; la ligne des apsides menée vers le périhélie; la perpendiculaire à celle-ci dans le plan de l’orbite; l’intersection du plan des \(x,y\) avec le plan de l’orbite. Nous dénoterons les points de la surface sphérique, auxquels répondent ces directions, par \(X,Y,Z,L,\Lambda,M,P,\Pi,N\). Ainsi \(X,Y,Z\) seront les pôles des grands cercles qui représentent les trois plans fondamentaux, et nous supposeront ces pôles pris du côté où les coordonnées sont censées positives; \(L\) sera le lieu héliocentrique de la planète et \(\Lambda\) le point du grand cercle de l’orbite avancé davantage que \(L\) de \(90^{0}\); \(P\) le périhélie; \(\Pi\) ce que devient \(\Lambda\) lorsque \(L\) devient \(P\); \(N\) sera le noeud ascendant de l’orbite sur le grand cercle dont \(Z\) est le pôle; enfin \(M\) sera l’un des pôles de l’orbite, et pour ne laisser aucune ambiguïté, nous supposerons que les points \(P,\Pi,M\) soient rangés sur la sphère dans le même ordre que les points \(X,Y,Z\). On voit que pour cet effet il faudra choisir \(M\) du même côté du grand cercle \(XY\) que \(Z\), ou du côté opposé, selon que le mouvement héliocentrique de la planète projeté sur le grand cercle \(XY\) va dans le sens de \(X\) vers \(Y\), ou dans le sens opposé: on nommera le mouvement \textit{direct} dans le premier cas, et \textit{rétrograde} dans le second. Pour distinguer analytiquement ces deux cas, soit \(\beta\) la distance de \(L\) au grand cercle \(XY\), ou plutôt \(\beta=90^{0}-LZ\), \(\lambda\) la longitude de la projection de \(L\) sur ce grand cercle comptée de \(X\) dans le sens de \(X\) vers \(Y\); nous aurons
\[
\begin{aligned}
x&=r\cos\beta\cos\lambda\\
y&=r\cos\beta\sin\lambda\\
z&=r\sin\beta.
\end{aligned}
\]
\par De là on tire \(x\mathrm{d}y-y\mathrm{d}x=rr\cos\beta^2\mathbin{.}\mathrm{d}\lambda\); ainsi dans le cas du mouvement \SourceLeaf{452}direct \(C\) sera positif, et négatif dans le cas du mouvement rétrograde: donc \(C\) et \(\cos ZM\) seront toujours du même signe. Maintenant on a
\[
\begin{array}{rl}
27)\quad&x=r\cos XL\\[4pt]
28)\quad&y=r\cos YL\\[4pt]
29)\quad&z=r\cos ZL
\end{array}
\]
et par un théorème connu
\[
\begin{aligned}
&\cos XL\mathbin{.}\cos XM+\cos YL\mathbin{.}\cos YM+\cos ZL\mathbin{.}\cos ZM\\
&\qquad=\cos LM=0.
\end{aligned}
\]
\par Par conséquent il sera
\[
x\cos XM+y\cos YM+z\cos ZM=0
\]
et puisque évidemment cette équation doit être identique à l’équation 7, les quantités \(A,B,C\) doivent être proportionnelles à \(\cos XM,\cos YM,\cos ZM\). Or par un autre théorème connu on a \(\cos XM^2+\cos YM^2+\cos ZM^2=1\); donc puisque nous venons de faire voir que \(\cos ZM\) et \(C\) sont nécessairement du même signe, et qu’il est de plus \(AA+BB+CC=mp\), nous devons conclure
\[
\begin{array}{rl}
30)\quad&A=\cos XM\mathbin{.}\BareRoot{mp}\\[4pt]
31)\quad&B=\cos YM\mathbin{.}\BareRoot{mp}\\[4pt]
32)\quad&C=\cos ZM\mathbin{.}\BareRoot{mp}.
\end{array}
\]
